% Options for packages loaded elsewhere
\PassOptionsToPackage{unicode}{hyperref}
\PassOptionsToPackage{hyphens}{url}
\documentclass[
  ignorenonframetext,
]{beamer}
\newif\ifbibliography
\usepackage{pgfpages}
\setbeamertemplate{caption}[numbered]
\setbeamertemplate{caption label separator}{: }
\setbeamercolor{caption name}{fg=normal text.fg}
\beamertemplatenavigationsymbolsempty
% remove section numbering
\setbeamertemplate{part page}{
  \centering
  \begin{beamercolorbox}[sep=16pt,center]{part title}
    \usebeamerfont{part title}\insertpart\par
  \end{beamercolorbox}
}
\setbeamertemplate{section page}{
  \centering
  \begin{beamercolorbox}[sep=12pt,center]{section title}
    \usebeamerfont{section title}\insertsection\par
  \end{beamercolorbox}
}
\setbeamertemplate{subsection page}{
  \centering
  \begin{beamercolorbox}[sep=8pt,center]{subsection title}
    \usebeamerfont{subsection title}\insertsubsection\par
  \end{beamercolorbox}
}
% Prevent slide breaks in the middle of a paragraph
\widowpenalties 1 10000
\raggedbottom
\AtBeginPart{
  \frame{\partpage}
}
\AtBeginSection{
  \ifbibliography
  \else
    \frame{\sectionpage}
  \fi
}
\AtBeginSubsection{
  \frame{\subsectionpage}
}
\usepackage{iftex}
\ifPDFTeX
  \usepackage[T1]{fontenc}
  \usepackage[utf8]{inputenc}
  \usepackage{textcomp} % provide euro and other symbols
\else % if luatex or xetex
  \usepackage{unicode-math} % this also loads fontspec
  \defaultfontfeatures{Scale=MatchLowercase}
  \defaultfontfeatures[\rmfamily]{Ligatures=TeX,Scale=1}
\fi
\usepackage{lmodern}
\ifPDFTeX\else
  % xetex/luatex font selection
\fi
% Use upquote if available, for straight quotes in verbatim environments
\IfFileExists{upquote.sty}{\usepackage{upquote}}{}
\IfFileExists{microtype.sty}{% use microtype if available
  \usepackage[]{microtype}
  \UseMicrotypeSet[protrusion]{basicmath} % disable protrusion for tt fonts
}{}
\makeatletter
\@ifundefined{KOMAClassName}{% if non-KOMA class
  \IfFileExists{parskip.sty}{%
    \usepackage{parskip}
  }{% else
    \setlength{\parindent}{0pt}
    \setlength{\parskip}{6pt plus 2pt minus 1pt}}
}{% if KOMA class
  \KOMAoptions{parskip=half}}
\makeatother
\usepackage{longtable,booktabs,array}
\usepackage{caption}
\captionsetup[table]{skip=6pt}
\newcounter{none} % for unnumbered tables
\usepackage{calc} % for calculating minipage widths
\usepackage{caption}
% Make caption package work with longtable
\makeatletter
\def\fnum@table{\tablename~\thetable}
\makeatother
\setlength{\emergencystretch}{3em} % prevent overfull lines
\providecommand{\tightlist}{%
  \setlength{\itemsep}{0pt}\setlength{\parskip}{0pt}}
% Manuscript macro/package fallbacks (newcommand -> providecommand).
% Core mathematical packages
\usepackage{amsmath,amssymb,amsthm}
\usepackage{mathtools}
% Keep long deployment prose and references within the text block.
% Algorithm formatting
% Tables
\usepackage{booktabs}
\usepackage{multirow}
% Graphics and floats
% Typography and layout
% Page layout (slightly smaller margins for academic journal style)
% Running headers/footers
% Caption formatting
% Colored boxes for callouts
% Cross-referencing with smart naming
% Hyperlink styling
% Configure cleveref for custom environments
% Theorem environments with consistent numbering
% Code listings and raw implementation examples in the manuscript
\usepackage{listings}
\lstset{basicstyle=\ttfamily\small,breaklines=true,columns=fullflexible}
% Math operators
\DeclareMathOperator*{\argmax}{arg\,max}
\DeclareMathOperator*{\argmin}{arg\,min}
\DeclareMathOperator{\sign}{sign}
\DeclareMathOperator{\KL}{KL}
\DeclareMathOperator{\Tr}{Tr}
% Custom commands for notation consistency
\providecommand{\calA}{\mathcal{A}}
\providecommand{\calB}{\mathcal{B}}
\providecommand{\calC}{\mathcal{C}}
\providecommand{\calD}{\mathcal{D}}
\providecommand{\calF}{\mathcal{F}}
\providecommand{\calG}{\mathcal{G}}
\providecommand{\calH}{\mathcal{H}}
\providecommand{\calI}{\mathcal{I}}
\providecommand{\calM}{\mathcal{M}}
\providecommand{\calO}{\mathcal{O}}
\providecommand{\calP}{\mathcal{P}}
\providecommand{\calS}{\mathcal{S}}
\providecommand{\calT}{\mathcal{T}}
\providecommand{\calW}{\mathcal{W}}
% Note: \Phi is a standard LaTeX command, do not redefine
% Trust notation
\providecommand{\trust}[2]{\mathcal{T}_{#1 \to #2}}
\providecommand{\trustt}[3]{\mathcal{T}_{#1 \to #2}^{#3}}
% Belief notation
\providecommand{\belief}[2]{\mathcal{B}_{#1}(#2)}
\providecommand{\belieft}[3]{\mathcal{B}_{#1}^{#2}(#3)}
% Cognitive state
\providecommand{\cogstate}[1]{\sigma_{#1}}
% Attack notation
\providecommand{\attack}[1]{\mathcal{A}_{#1}}
\providecommand{\adversary}[1]{\Omega_{#1}}
% Defense notation
\providecommand{\firewall}{\mathcal{F}}
\providecommand{\sandbox}{\mathcal{S}_{box}}
% Probability and expectation
\providecommand{\E}{\mathbb{E}}
\providecommand{\Prob}{\mathbb{P}}
\providecommand{\indicator}{\mathbb{1}}
% QED symbol
\providecommand{\qedsymbol}{$\blacksquare$}
% Front matter opens on the cover/title page (no list of figures/tables precedes it).

% Slide layout defaults for warning-clean scientific decks.
\usepackage{etoolbox}
\IfFileExists{xurl.sty}{\usepackage{xurl}}{}
\IfFileExists{seqsplit.sty}{\usepackage{seqsplit}}{\newcommand{\seqsplit}[1]{#1}}
\protected\def\breakseq#1{\seqsplit{#1}}
\protected\def\breaktt#1{\begingroup\ttfamily\seqsplit{#1}\endgroup}
\setlength{\emergencystretch}{6em}
\tolerance=5000
\hbadness=10000
\hfuzz=1pt
\setlength{\tabcolsep}{2pt}
\AtBeginEnvironment{longtable}{\tiny\renewcommand{\arraystretch}{0.86}\setlength{\tabcolsep}{1pt}}
\AtBeginEnvironment{tabular}{\tiny\renewcommand{\arraystretch}{0.86}\setlength{\tabcolsep}{1pt}}
\AtBeginEnvironment{equation}{\tiny}
\AtBeginEnvironment{equation*}{\tiny}
\AtBeginEnvironment{align}{\tiny}
\AtBeginEnvironment{align*}{\tiny}
\AtBeginEnvironment{itemize}{\footnotesize}
\AtBeginEnvironment{enumerate}{\footnotesize}
\AtBeginEnvironment{description}{\footnotesize}
\setbeamerfont{caption}{size=\tiny}
\setbeamerfont{caption name}{size=\tiny}
\setbeamerfont{normal text}{size=\small}
\setbeamerfont{frametitle}{size=\small}
\setbeamerfont{section title}{size=\footnotesize}
\setbeamerfont{subsection title}{size=\footnotesize}
\setbeamertemplate{section page}{%
  \centering
  \begin{beamercolorbox}[sep=12pt,center,wd=\paperwidth]{section title}
    \parbox{0.86\paperwidth}{\centering\usebeamerfont{section title}\insertsection\par}
  \end{beamercolorbox}
}
\setbeamertemplate{subsection page}{%
  \centering
  \begin{beamercolorbox}[sep=8pt,center,wd=\paperwidth]{subsection title}
    \parbox{0.86\paperwidth}{\centering\usebeamerfont{subsection title}\insertsubsection\par}
  \end{beamercolorbox}
}
\setlength{\abovecaptionskip}{2pt}
\setlength{\belowcaptionskip}{0pt}

% Natbib and cross-reference fallbacks — slides don't load natbib
% or cleveref, but combined-PDF manuscript prose may emit these
% commands. The fallback renders citations as a bracketed key list
% and cross-references as detokenized labels so slides don't fail on
% undefined control sequences or raw underscores. \providecommand is
% a no-op if packages load later.
\providecommand{\citep}[1]{[#1]}
\providecommand{\citet}[1]{#1}
\providecommand{\citealp}[1]{#1}
\providecommand{\citeauthor}[1]{#1}
\providecommand{\citeyear}[1]{#1}
\providecommand{\cref}[1]{\texttt{\detokenize{#1}}}
\providecommand{\Cref}[1]{\texttt{\detokenize{#1}}}
\providecommand{\crefrange}[2]{\texttt{\detokenize{#1}}--\texttt{\detokenize{#2}}}
\providecommand{\Crefrange}[2]{\texttt{\detokenize{#1}}--\texttt{\detokenize{#2}}}
\providecommand{\paragraph}[1]{\textbf{#1}\ }

% Environment-providing packages carried over from the manuscript.
\usepackage{algpseudocode}

% Non-floating stand-in for the `algorithm` float.
\newenvironment{algorithm}[1][]{%
  \par\medskip\noindent\rule{\linewidth}{0.4pt}\par\nobreak\small
  \renewcommand{\caption}[1]{\par\noindent\textbf{##1}\par}}{%
  \par\nobreak\noindent\rule{\linewidth}{0.4pt}\par\medskip}

\newtheorem{proposition}{Proposition}
\newtheorem{hypothesis}{Hypothesis}
\newtheorem{remark}[theorem]{Remark}
\newtheorem{axiom}{Axiom}
\newtheorem{property}{Property}
\usepackage{bookmark}
\IfFileExists{xurl.sty}{\usepackage{xurl}}{} % add URL line breaks if available
\urlstyle{same}
% fallback for those not using the hyperref driver hyperxmp:
\makeatletter
\@ifundefined{xmpquote}{\newcommand{\xmpquote}[1]{#1}}{}
\makeatother
\hypersetup{
  hidelinks,
  pdfcreator={LaTeX via pandoc}}

\author{\texorpdfstring{}{}}
\date{}

\begin{document}

\begin{frame}
\newpage
\end{frame}

\section{Why Cognitive Security Matters Now}\label{sec:introduction}

\begin{frame}[allowframebreaks]{The Operational Reality}
\protect\phantomsection\label{the-operational-reality}
Something fundamental changed in how AI systems work, and the security
community is catching up.

In 2023, AI security often meant preventing chatbots from saying things
they shouldn't. The attack surface was a text box; the defense was a
filter.

By 2026, we are securing \textbf{multiagent operators}---networks of
specialized AI agents that delegate to each other, form beliefs about
each other's outputs, build trust relationships over time, and take
actions with real-world consequences. These systems write code, manage
infrastructure, and move money.

\end{frame}

\begin{frame}[allowframebreaks]{The Operational Reality}
The shift is from ``content safety'' to ``cognitive integrity.'' The
risk isn't just that an agent says something wrong, but that it
\emph{believes} something wrong---and acts on it.
\end{frame}

\begin{frame}[allowframebreaks]{The Good News: It's Solvable}
\protect\phantomsection\label{the-good-news-its-solvable}
This is not a theoretical warning about future doom. It is an
engineering problem with established solutions.

The Cognitive Integrity Framework (CIF) was developed to secure these
systems, and the companion papers in this series demonstrate its
efficacy from formal, computational, and applied angles.

\end{frame}

\begin{frame}[allowframebreaks]{The Good News: It's Solvable}
\begin{itemize}
\tightlist
\item
  \textbf{Part 1: Formal Foundations} (DOI: 10.5281/zenodo.18364119)
  proved that trust can be mathematically bounded. We defined the
  ``Trust Calculus'' which guarantees that no matter how clever an
  adversary is, they cannot amplify their influence through delegation
  chains. It also introduces the Defense Composition Algebra, the
  five-tier adversary taxonomy (\(\Omega_1\)--\(\Omega_5\)), and
  information-theoretic stealth-impact bounds.
\item
  \textbf{Part 2: Computational Validation} (DOI:
  10.5281/zenodo.18364128) implemented this theory in Python and tested
  it against a corpus of 950 attacks across four production
  architectures, reporting ablation studies, Bayesian uncertainty
  quantification, colony-scale benchmarks at 20--100 agents, and a
  category-theoretic formalization of defense composition (Defense
  Category \(\calD\), Theorems CT.1--CT.3) with a composable
  visualization engine and interactive CIF Composer web UI.
\item
  \textbf{Applications (§9--§10, this paper):} The integrated
  CIF-AD-OODA analytical model is applied across ten critical domains
  (rare-earth mining, nation-state alliances, cyber-security, drone
  warfare, supply chain, biowarfare, food security, trade wars,
  infrastructure, information ecosystems), identifying three universal
  attack patterns and four novel defense extensions.
\end{itemize}

\end{frame}

\begin{frame}[allowframebreaks]{The Good News: It's Solvable}
The combined evidence includes \textbf{3,390 tests} and a
\textbf{96--100\% parametric detection ceiling} across all attack
categories and architectures (Part 2), alongside a lower real-pipeline
multi-seed mean of approximately 44.8\% (30 seeds). Direct-injection
detection reaches 99--100\% in the fully defended parametric
configuration; plus CIF coverage is analyzed across all ten operational
domains in §9--§10 with retrospective analysis of six documented
2024--2025 AI-agent incidents.
\end{frame}

\begin{frame}[allowframebreaks]{The Purpose of This Guide}
\protect\phantomsection\label{the-purpose-of-this-guide}
We wrote Part 1 for the theorists, Part 2 for the experimentalists, and
the Applications section (§9--§10) for domain experts. We wrote the
Practitioner section (§1--§8)---the opening half of this unified
paper---to translate those findings into deployable engineering
practice.

Our goal is to describe how the defenses validated in the companion
papers can be architected in production systems. We focus on the
practical application of the formal proofs:

\end{frame}

\begin{frame}[allowframebreaks]{The Purpose of This Guide}
\begin{itemize}
\tightlist
\item
  How the \textbf{Trust Decay} factor (\(\delta\)) functions in
  different topologies.
\item
  How \textbf{Behavioral Tripwires} served as effective detection
  mechanisms for hallucination.
\item
  How the \textbf{Cognitive Firewall} filtered inputs before they became
  beliefs.
\end{itemize}
\end{frame}

\begin{frame}[allowframebreaks]{How to Use This Resource}
\protect\phantomsection\label{how-to-use-this-resource}
\end{frame}

\begin{frame}[allowframebreaks]{How to Use This Resource}
\begin{itemize}
\tightlist
\item
  \textbf{Section 2} summarizes the theoretical concepts from Part 1,
  providing the necessary vocabulary.
\item
  \textbf{Section 3} reviews the empirical evidence from Part 2,
  detailing which architectures performed best against specific threats.
\item
  \textbf{Section 4} analyzes the attack scenarios used in our testing
  corpus. For domain-specific attack patterns (FR Polarity Inversion,
  Constraint Relaxation, Context Boundary Violation) and documented
  real-world AI-agent incidents, see \textbf{§9--§10} (the Applications
  section of this paper).
\item
  \textbf{Section 5} presents the specific configuration profiles that
  yielded the highest security margins in simulation, with deployment
  guides, incident response playbooks, monitoring, and cost--benefit
  analysis.
\item
  \textbf{Sections 6-7} discuss the limitations discovered during
  testing and the open problems that remain; domain-specific case
  studies and novel defense extensions (verification channel separation,
  active perturbation probing, physics-informed invariants) are treated
  in \textbf{§9--§10}.
\end{itemize}

\end{frame}

\begin{frame}[allowframebreaks]{How to Use This Resource}
This paper serves as a report on the current state of cognitive security
engineering, grounded in the data and definitions of the CIF series.
Readers seeking derivations or proofs should consult Part 1; readers
seeking empirical measurements should consult Part 2; readers evaluating
CIF for a specific operational sector should consult §9--§10 of this
paper.
\end{frame}

\begin{frame}[allowframebreaks]{Reading Companion: Where to Find
Specific Topics}
\protect\phantomsection\label{sec:reading-companion}
This paper is designed to stand alone as the practitioner's reference of
the series. Where a concept or technique is developed more fully
elsewhere, the table below points the way.

\end{frame}

\begin{frame}[allowframebreaks]{Reading Companion: Where to Find
Specific Topics}
{\def\LTcaptype{none} % do not increment counter
\begin{longtable}[]{@{}
  >{\raggedright\arraybackslash}p{(\linewidth - 2\tabcolsep) * \real{0.5714}}
  >{\raggedright\arraybackslash}p{(\linewidth - 2\tabcolsep) * \real{0.4286}}@{}}
\toprule\noalign{}
\begin{minipage}[b]{\linewidth}\raggedright
If you want\ldots{}
\end{minipage} & \begin{minipage}[b]{\linewidth}\raggedright
\ldots consult\ldots{}
\end{minipage} \\
\midrule\noalign{}
\endhead
Formal definitions, proofs, and theorems (Trust Calculus, Defense
Composition Algebra, stealth--impact bound) & \textbf{Part 1} (DOI:
10.5281/zenodo.18364119), §§4--5, 7 \\
Adversary taxonomy \(\Omega_1\)--\(\Omega_5\) formal characterization &
\textbf{Part 1}, its Threat Model section (Formal Characterization of
the Adversary Classes) \\
Model-checked safety invariants + NuSMV/TLA+ specifications &
\textbf{Part 1}, ``Formal Verification: Safety Properties and Model
Checking'', \textbf{Part 2} S04 \\
Eusocial-colony analogy (biological existence proof for CIF-like
architectures) & \textbf{Part 1} S02 \\
950-attack corpus generation, examples, ethics & \textbf{Part 2} (DOI:
10.5281/zenodo.18364128), §3 + S03 \\
Detailed detection rates per architecture (Claude Code, AutoGPT, CrewAI,
LangGraph) & \textbf{Part 2}, ``Extended Experimental Results'' for
measured Claude Code and CrewAI rates, and ``Per-Architecture Parametric
Detection Rates'' plus ``Cross-Architecture Parametric Summary'' in S08
for all four architectures \\
Ablation studies + Bayesian uncertainty & \textbf{Part 2}, ``Ablation
Studies and Scalability Benchmarks'' and ``Bayesian Uncertainty
Quantification'' \\
Parametric design-level ceiling (96--100\%) & \textbf{Part 2} S08 \\
Game-theoretic adversarial analysis / Nash equilibrium & \textbf{Part
2}, the ``Game-Theoretic Analysis'' subsection of ``Theoretical
Connections'' for the payoff matrix and Theorem GT.1, and
``Game-Theoretic Arms Race Dynamics'' in the Discussion \\
Category-theoretic formalization of defense composition (Defense
Category \(\calD\), Theorems CT.1--CT.3) & \textbf{Part 2}, ``Defense
Composition as Category Theory'' in ``Theoretical Connections'' for CT.1
and CT.2, and ``Composability Algebra: Monadic Defense Chains'' for
CT.3, with the extended treatment in ``Category-Theoretic Foundations of
Defense Composition'' \\
Composable visualization engine + CIF Composer interactive web UI &
\textbf{Part 2} (output/web/cif\_composer.html) \\
Free-energy connections (FEP.1--FEP.2) & \textbf{Part 2} Theoretical
Connections (Active Inference and the Free Energy Principle), Supplement
S10 Information Geometry of Belief Manipulation \\
Framework API reference + pseudocode & \textbf{Part 2} S05, S07 \\
Application of CIF to specific operational sectors (10 domains analyzed)
& \textbf{§9--§10} (this paper) \\
Three universal attack patterns across domains (FR Polarity Inversion,
Constraint Relaxation, Context Boundary Violation) & \textbf{§10} (this
paper) \\
Four novel defense extensions (verification channel separation, active
perturbation probing, physics-informed invariants, semiotic decoupling)
& \textbf{§9} (this paper) \\
Retrospective mapping of 2024--2025 AI-agent security incidents (Replit,
Copilot RCE, Slack AI, \$3.2M procurement fraud, etc.) & \textbf{S3}
(this paper) \\
CIF-AD-OODA integration model for goal-hijacking & \textbf{§9} (this
paper) \\
\bottomrule\noalign{}
\end{longtable}
}

\end{frame}

\begin{frame}[allowframebreaks]{Reading Companion: Where to Find
Specific Topics}
\textbf{Code and Repository}: The companion codebase, attack corpus
documentation, and deployment tooling are maintained at
\url{https://github.com/docxology/cognitive_integrity} (DOI:
10.5281/zenodo.18364130; companion parts: Part 1 DOI
10.5281/zenodo.18364119, Part 2 DOI 10.5281/zenodo.18364128).
\end{frame}

\end{document}
